\documentclass[12pt,a4paper]{article}
\usepackage[utf8]{inputenc}
\usepackage[T1]{fontenc}
\usepackage{geometry}
\usepackage{graphicx}
\usepackage{booktabs}
\usepackage{float}
\usepackage{xcolor}
\usepackage{amsmath}
\usepackage{listings}
\geometry{margin=2.4cm}
\graphicspath{{./}}
\renewcommand{\arraystretch}{1.2}
\definecolor{codebg}{RGB}{246,248,250}
\renewcommand{\figurename}{Figura}
\renewcommand{\tablename}{Tabela}
\lstset{basicstyle=\ttfamily\small,backgroundcolor=\color{codebg},frame=single,
  breaklines=true,columns=fullflexible,showstringspaces=false}

\title{Análise da Utilização do Sistema por Estação do Ano}
\author{}
\date{}

\begin{document}
\maketitle

\section{Objetivo e critério de baixa utilização}

Esta análise compara a variável \texttt{total\_user} entre as quatro estações
do ano. Para cada estação, foram calculados a média, a mediana, o desvio padrão
amostral e a proporção de dias classificados como \texttt{low\_usage}.

Adotou-se a definição da Questão 2:
\[
\texttt{low\_usage}=\begin{cases}
\texttt{TRUE}, & \texttt{total\_user}\leq Q_1,\\
\texttt{FALSE},& \texttt{total\_user}>Q_1.
\end{cases}
\]

Antes do cálculo, os 300 valores de \texttt{total\_user} foram organizados em
ordem crescente. No R, foi usado \texttt{sort()} explicitamente e o método
\texttt{type = 7}, padrão de \texttt{quantile()} e compatível com o cálculo dos
quartis empregado pelo \texttt{boxplot()}.

Para $n=300$, as posições teóricas do primeiro e terceiro quartis pelo método 7
são
\[
h_{Q_1}=1+(300-1)\times0{,}25=75{,}75,
\qquad
h_{Q_3}=1+(300-1)\times0{,}75=225{,}25.
\]
Consequentemente, há interpolação entre os valores ordenados adjacentes. Foram
obtidos $Q_1=2105{,}5$, $Q_2=3996$ e $Q_3=4681$. Assim, 75 dias foram
classificados como baixa utilização.

\begin{figure}[H]
\centering
\includegraphics[width=.92\textwidth]{print_quartis.png}
\caption{Saída dos cálculos dos quartis e da classificação \texttt{low\_usage}.}
\end{figure}

\section{Código utilizado no R}

\begin{lstlisting}[language=R]
library(readxl)
library(dplyr)
library(ggplot2)

dados <- read_excel("2.1.total_user.xlsx")

quartis <- quantile(sort(dados$total_user),
                    probs = c(0.25, 0.50, 0.75), type = 7)
Q1 <- quartis[[1]]

dados <- dados |>
  mutate(low_usage = total_user <= Q1)

resumo_estacoes <- dados |>
  group_by(season) |>
  summarise(
    n = n(),
    media = mean(total_user),
    mediana = median(total_user),
    desvio_padrao = sd(total_user),
    dias_low_usage = sum(low_usage),
    proporcao = mean(low_usage),
    .groups = "drop"
  )
print(resumo_estacoes)
\end{lstlisting}

\section{Resultados por estação}

\begin{table}[H]
\centering
\small
\caption{Medidas descritivas e baixa utilização por estação.}
\begin{tabular}{lrrrrrr}
\toprule
Estação & $n$ & Média & Mediana & Desvio padrão & Dias low & Proporção \\
\midrule
Inverno    & 73 & 1.639,82 & 1.605,0 & 575,69   & 60 & 82,19\% \\
Primavera  & 92 & 3.775,17 & 4.085,5 & 1.138,90 & 11 & 11,96\% \\
Verão      & 94 & 4.464,36 & 4.615,5 & 798,35   & 3  & 3,19\% \\
Outono     & 41 & 4.080,27 & 4.304,0 & 1.054,78 & 1  & 2,44\% \\
\bottomrule
\end{tabular}
\end{table}

\begin{figure}[H]
\centering
\includegraphics[width=\textwidth]{print_resumo_estacoes.png}
\caption{Saída do resumo por estação. A proporção é apresentada em escala de 0 a 1.}
\end{figure}

O inverno apresenta a menor utilização: média de 1.639,82 e mediana de 1.605
usuários. Além disso, 60 dos 73 dias de inverno (82,19\%) estão no primeiro
quartil global. Na primavera, a média sobe para 3.775,17 e a proporção de baixa
utilização cai para 11,96\%.

O verão possui a maior média (4.464,36) e a maior mediana (4.615,5), com apenas
3,19\% dos dias classificados como \texttt{low\_usage}. O outono apresenta a
segunda maior mediana (4.304) e a menor proporção de baixa utilização (2,44\%).
Em termos absolutos, a primavera tem o maior desvio padrão; portanto, foi a
estação com maior dispersão absoluta dos totais diários.

\section{Boxplot por estação}

\begin{lstlisting}[language=R]
ggplot(dados, aes(x = factor(season,
  levels = unique(dados$season)),
  y = total_user, fill = season)) +
  geom_boxplot(width = 0.65, alpha = 0.80) +
  geom_hline(yintercept = Q1, linetype = "dashed",
             colour = "gray35") +
  labs(title = "Distribuicao do total de usuarios por estacao",
       x = "Estacao do ano", y = "Total diario de usuarios") +
  guides(fill = "none") +
  theme_minimal(base_size = 12)
\end{lstlisting}

\begin{figure}[H]
\centering
\includegraphics[width=.94\textwidth]{boxplot_estacoes.png}
\caption{Boxplot do total diário de usuários por estação. A linha tracejada
indica o limite global de baixa utilização, $Q_1=2105{,}5$.}
\end{figure}

O boxplot confirma a separação do inverno em relação às demais estações: sua
caixa e sua mediana estão abaixo das outras, e o terceiro quartil do inverno
($1969$) também fica abaixo do limite global de baixa utilização. Primavera e
outono têm caixas mais amplas, coerentes com seus desvios padrão elevados. O
verão concentra valores em um patamar superior e apresenta distribuição mais
compacta que primavera e outono.

Os pontos fora dos bigodes são observações atípicas segundo a regra de
$1{,}5\times IQR$; eles não são necessariamente erros e foram mantidos na
análise. O gráfico e as estatísticas mostram associação descritiva entre estação
e nível de utilização, mas não estabelecem causalidade.

\section{Conclusão}

Os dias de menor utilização concentram-se fortemente no inverno: essa estação
reúne 60 dos 75 dias classificados como \texttt{low\_usage}. Verão apresenta o
maior nível típico de utilização, enquanto outono registra a menor proporção de
dias abaixo ou iguais ao primeiro quartil global. Logo, dentro do período
observado, o inverno é a característica sazonal mais associada aos dias de
baixa utilização do sistema.

\end{document}

